%%%%%%%%%%%%%%%%%%%%%%%%%%%%%%%%%%%%%%%%%
% LaTeX Templates Journal Article
% LaTeX Class
% Version 2.0 (February 7, 2023)
%
% This class originates from:
% https://www.LaTeXTemplates.com
%
% Author:
% Vel (vel@latextemplates.com)
%
% License:
% CC BY-NC-SA 4.0 (https://creativecommons.org/licenses/by-nc-sa/4.0/)
%
%%%%%%%%%%%%%%%%%%%%%%%%%%%%%%%%%%%%%%%%%

\documentclass{ExcelAtFIT}
%\ExcelFinalCopy

\hypersetup{
	pdftitle={Paper Title},
	pdfauthor={Author}
  pdfkeywords={Keywords}
}

\lstset{ 
	backgroundcolor=\color{white},   % choose the background color; you must add \usepackage{color} or \usepackage{xcolor}; should come as last argument
	basicstyle=\footnotesize\tt,        % the size of the fonts that are used for the code
}

\usepackage{float}
\usepackage{hyphenat}

\addbibresource{references.bib}

\providecommand{\keywords}[1]
{
  \small	
  \textbf{\textit{Keywords---}} #1
}

\ExcelYear{2026}

\PaperTitle{(((title)))}

\Authors{Axiom K. Playwright (AI Author)*}
\affiliation{*%
  {The Ma'ayan Laboratory, Mount Sinai Center for Bioinformatics, Department of Pharmacological Sciences, Windreich Department of Artificial Intelligence and Human Health, Icahn School of Medicine at Mount Sinai, New York, NY 10029}}


\Abstract{
  (((abstract)))
}

\Keywords{CFDE, Gene Matrix Transpose, Gene Set, Enrichr, LINCS, GTEx, GlyGen, Aging, Glycosylation}

\begin{document}

\startdocument
\raggedbottom

\section*{Introduction}\label{Introduction}

Common Fund programs have generated many diverse and robust multiomics datasets in the past 20 years.
However, knowledge contained in these datasets can be limited if only considered within the context of any individual dataset.
To address this, we have converted Common Fund datasets from the Common Fund Data Ecosystem (CFDE) into gene set libraries  to combine multiple datasets at a time and investigate interesting overlapping features \autocite{CFDE}.
Pairwise gene set library crossing has been explored previously in GeneSetCart \autocite{GeneSetCart}, and Harmonizome 3.0 \autocite{Harmonizome}, where gene sets from two libraries where significant overlaps were found using Fisher's exact test.
These top crossings could then be used to generate LLM hypotheses by augmenting them with enriched biological terms. Here we extend the strategy by combining more libraries at once.
Each combination of three gene sets (one from each library) is analyzed to identify signficant pairwise intersections and a robust 3-way intersecting gene set.

For this report, we combined gene sets from the (((dataset_info.1.name))) \autocite{(((dataset_info.1.resource)))}, (((dataset_info.2.name))) \autocite{(((dataset_info.2.resource)))}, and (((dataset_info.3.name))) \autocite{(((dataset_info.3.resource)))} datasets.
Among the top-ranked crossings, we identified an overlap with rank (((rank))) and a p-value of (((p_value))) warranting further investigation.
(((introduction)))

\section*{Methods and Results}\label{Methods and results}

\subsection*{Dataset Processing}\label{Dataset Processing}

Gene matrix transpose (GMT) files used for crossing were created for the Common Fund Data Ecosystem (CFDE) Workbench Data Portal \autocite{CFDE}.

(((dataset_info.1.processing)))

(((dataset_info.2.processing)))

(((dataset_info.3.processing)))

\subsection*{Intersecting Genes}\label{Intersecting Genes}

\begin{figure}[h]
  \includegraphics[width=\columnwidth]{figures/venn\_diagram.pdf}
  \caption{A venn diagram of the crossed gene sets. Sizes of each set and 2-way and 3-way intersection are labelled.}
  \label{fig:venn}
\end{figure}

Crossing results were ranked by hypergeometric test p-values computed from pairwise overlaps for each collection of gene sets.
The intersection has a p-value of (((p_value))), a Jaccard index of (((jaccard))), and (((n_genes))) intesecting genes: (((intersecting_genes))).

\subsection*{Enrichment Results}\label{Enrichment Results}

The gene list derived from the crossing intersection was uploaded to Enrichr to explore functional enrichment \autocite{Enrichr}.
The submission used the default background and queried the GO Biological Process 2025 \autocite{GO}, KEGG Pathways 2026 \autocite{KEGG}, ChEA Transcription Factor Targets 2022 \autocite{ChEA}, and the GWAS Catalog Phenotypes 2025 \autocite{GWASCatalog} libraries.
This allowed the platform to return a comprehensive set of over-represented terms associated with the intersecting genes and further inform the biological context.

\begin{figure}[ht!]
  \includegraphics[width=\columnwidth]{figures/enrichment_bar.pdf}
  \caption{A bar plot off the enriched terms obtained by submitting the intersecting gene set to Enrichr from the \textbf{A.} GO Biological Processes 2025, \textbf{B.} KEGG Pathways 2026, \textbf{C.} ChEA Transcription Factor Targets 2022, and \textbf{D.} GWAS Catalog 2025 libraries}
  \label{fig:enrichment}
\end{figure}

\section*{Discussion}\label{Discussion}

(((gene_summaries)))

(((enrichment_results)))

(((discussion)))

\section*{Acknowledgements}\label{Acknowledgements}
This project was funded by the NIH grant \href{https://reporter.nih.gov/search/SdeFoZSP2U2zRTjMZKFHlQ/project-details/11080094}{OT2OD036435}.
The OpenAI gpt-oss-120b LLM was used to assist in writing the article.

\printbibliography

\end{document}